

\section[Scalable]{Scalabilité \& Tolérance aux pannes}

\subsection{Distribution \& Réplication}

%%%
\begin{frame}[fragile]{Distribution}

Cassandra s'appuie sur le hachage cohérent, très inspiré du système Dynamo (Amazon)

\begin{itemize}
	\item Anneau sur l'espace $[-2^{63}, 2^{63}]$
	  
	  
	  \item Les serveurs ont une position (\textit{token}) calculée ou attribuée
	      explicitement.
	      
	    \item Le découpage peut être lissé grâce à un système de nœuds virtuels (256 par nœud physique).
\end{itemize}

\vfill

Des serveurs peuvent être ajoutés ou supprimés à tout moment.

\end{frame}

%%%
\begin{frame}[fragile]{Routage des requêtes}

Chaque nœud connaît la topologie complète de l'anneau.
    
\begin{itemize}
          \item Chaque mise à jour (ajout suppression de serveur) est propagée à l'ensemble
              des participants.
          \item La table peut devenir assez volumineuse (surtout avec nœuds virtuels): 
\end{itemize}
      
\vfill

Cassandra fonctionne en \alert{multi-nœuds}.

\begin{itemize}
          \item Un requête peut être prise en charge par n'importe quel nœud, le \alert{coordinateur}
          \item Le coordinateur identifie les nœuds responsables de la clé du nouveau document
\end{itemize}

\end{frame}

%%%
\begin{frame}[fragile]{Réplication}

La réplication est fixée au niveau du \textit{keyspace}: 3 par défaut.
	
\vfill

On indique également une \textit{replication strategy}.

\begin{minted}{sql}
     CREATE keyspace repli 
              with replication = {'class':'SimpleStrategy', 
                                 'replication_factor':3};
\end{minted}

\vfill

La stratégie de réplication "simple" copie sur les successeurs du nœud-cible de l'insertion.

\vfill

Une stratégie plus complexe tient compte de la topologie du réseau (niveaux \textit{racks}
et \textit{data centers}). 

\end{frame}


\begin{frame}[fragile]{Ecriture et cohérence des données}

Cassandra propose un mécanisme avancé de gestion du compromis entre cohérence et latence. 


\vfill

\figSlideWithSize{cassandra-log}{10cm}



\end{frame}

\begin{frame}[fragile]{Paramétrage de la cohérence en écriture}

Le niveau indique le nombre d'acquittements que le coordinateur doit rececoir des nœuds de stockage avant d'acquitter
à son tour le client.

\begin{itemize}
   \item \texttt{ANY}:  on accepte 0 acquittements! Le document est conservé  dans une zone temporaire, en attente d'une nouvelle tentative d'écriture...

    \item \texttt{ONE} : Au moins un acquittement

    \item \texttt{QUORUM} : Au moins $\lfloor replication/2 \rfloor +1$ acquittements, soit 2 pour un facteur de réplication 3.
    
    \item \texttt{ALL} : La réponse au client sera assurée lorsque la ressource aura été écrite dans tous les réplicas. 
\end{itemize}

\vfill

Rappel: plus on demande d'acquittements, plus on privilégie la cohérence sur la latence.

\vfill

Cassandra, système réputé très rapide en insertions.

\end{frame}

 
\begin{frame}[fragile]{Les stratégies en lecture}

Même type de paramétrage en lecture. La version avec l'estampille la plus récente est renvoyée au client.

\begin{itemize}
    \item \texttt{ONE} : le coordinateur transmet la première réponse reçue.

    \item \texttt{QUORUM} : Le coordinateur reçoit la réponse de au moins $\lfloor replication/2 \rfloor +1$  
       réplicas, et renvoie au client la ressource avec l'estampille la plus récente.
    
    \item \texttt{ALL} : Le coordinateur attend d'avoir reçu la réponse  de tous les réplicas. 
\end{itemize}

\vfill

Même remarque: compromis latence/cohérence
\end{frame}
 

\begin{frame}[fragile]{Comment régler les paramètres?}

Soit
    \begin{itemize}
        \item \textbf{N} : Taux de réplication
        \item \textbf{W} : Nb minimum d'écritures devant accuser de réception
        \item \textbf{R} : Nb copies d'une donnée à consulter pour une requête
    \end{itemize} 

\alert{Alors, si $ W + R > N$, on assure la cohérence des lectures.}

\vfill

Exemple

\begin{itemize}
    \item Soit N=3, W=ALL, R=1: toute lecture lit un document à jour
    \item Soit N=3, W=1, R=ALL: une des lectures lit le document à jour, et c'est celui-ci qui est renvoyé
    \item Soit N=3, W=QUORUM, R=QUORUM: une des lectures lit forcément lune des versions  à jour, et c'est celui-ci qui est renvoyé
    
\end{itemize}

\vfill

Simple et élégant!

\end{frame}

\end{document}

\begin{frame}[fragile]{Cohérence : Options de Cohérence}
\begin{itemize}
	\item 	\begin{minted}{sql}
	CONSISTENCY <level> ;
	\end{minted}
	\textit{level} : \footnotesize \textsc{ANY, ONE, TWO, THREE, QUORUM, ALL}\normalsize
	\item \textbf{\textsc{ONE/ANY}} : par défaut, au moins 1 acquittement	\footnotesize (Bon compromis)\normalsize
	\item \textbf{\textsc{ALL}} : W tous acquittés, 1 seul R \footnotesize (Système avec peu d'écritures)\normalsize
	\item \textbf{\textsc{QUORUM}} : N/2 + 1 (pour R et W) \\\footnotesize (Véritable quorum, possibilité d'incohérence - faible)\normalsize
	
\end{itemize}
\end{frame}